\title{Petrol Pal}

\maketitle
\section*{Project Description}
PetrolPal is a platform that makes it easy for you to find the best petrol prices. 
It lets you search for and view petrol prices at nearby stations, helping you make informed decisions and save on petrol.

\section*{Project Deliverables}
An MVP will require the following functionalities:
\begin{enumerate}
\item An interactive spatial visualisation allowing users to view real-time fuel pricing data across various geographic coordinates via a browser-based map. 
\item Integration of geolocation services to accurately render the user's current physical position within the map interface. 
\item A dynamic filtering engine that enables users to refine map results based on categorical data, such as fuel type (e.g., E10, Unleaded 91, Diesel, or Premium). 
\item A conditional authentication system where public data viewing is accessible without an account, while contribution and notification features require secure user registration. 
\item A crowdsourced submission mechanism for price updates, requiring the upload of photo evidence to cloud storage for verification purposes. 
\item A notification system that monitors user-defined "favourite" geographic zones and triggers alerts based on price thresholds or market trends. 
\item Spatial query functionality allows users to interact with specific map markers to retrieve detailed metadata and current pricing for a specific station.
\end{enumerate}

\section*{Quality Attributes}
Below are the main quality attributes the application will focus on. 
In addition to these, implement strong privacy and security measures to protect user data.
\begin{enumerate}
\item Availability: PetrolPal must be accessible via a mobile-responsive browser or application at all times to accommodate the urgent and 24/7 nature of fuel purchasing. The system will be evaluated on its ability to maintain constant uptime and provide up-to-date pricing data across all hours of operation. Basic tests must be created to monitor and ensure that the system performs under different environments and maintains accessibility during simulated 24/7 operational cycles.
\item Scalability: The platform's scalability should be evaluated by its ability to maintain performance metrics during unpredictable traffic spikes, such as those during holiday weekends or at the low points of a fuel price cycle. The system must demonstrate elasticity by dynamically adjusting resources to handle high-concurrency periods without impacting the user experience. Evaluation will be based on database performance under simulated load increases.
\item Reliability: The system is considered reliable if the pricing database maintains strict data consistency across multiple replicas, ensuring that all users see the most accurate and up-to-date information. Evaluation will involve testing the system’s ability to handle concurrent price updates from multiple users and demonstrating that data is synced correctly across the infrastructure. If a failure occurs, there must be a recovery strategy, such as instance snapshots, and a clear method to determine the most accurate data source in the event of a sync discrepancy.
\end{enumerate}

\end{document}